% This is a report section, not a standalone LaTeX document.
\section{A broader source associativity criterion}
\label{sec:broad-associativity}

The search for an additional state reduction led to a source transformation
with fewer hypotheses than the earlier closed-chunk rule.  It is an
independent source result; no smaller binary machine is inferred from it.

Write $\langle P,Q\rangle$ for the binary source constructor.  A local
instruction has an ordinary successor, an explicit break or continue
target, or a halt outcome.  Once an instruction has exhausted the zeros
or ones in its path inside a subtree, an escaping break or continue still
requires a positive number of the corresponding ancestors outside that
subtree.  Ordinary fallthrough is denoted $B_0$; $B_k$ for $k\geq1$ and
$C_k$ for $k\geq1$ denote escaping break and continue targets with the
indicated residual counts.

\paragraph{General associativity lemma.}
Let $A,B,C$ be nonempty source trees.  If $A$ has no escaping break and $C$
has no escaping continue, then
\[
 \langle A,\langle B,C\rangle\rangle
 \quad\longleftrightarrow\quad
 \langle\langle A,B\rangle,C\rangle
\]
preserves every ordered instruction and every possible control target in
every surrounding source context in which the required external ancestor
targets are defined.  The middle tree $B$ is arbitrary.
Halt instructions are permitted in all three trees.  Consequently, the
identity on ordered instruction positions and register files gives a
register-machine bisimulation, including divergent computations.

\paragraph{Proof.}
Let $a$ and $b$ be the positive leaf counts of $A$ and $B$.  Both trees
place their instruction lists at the same offsets $0,a,a+b$.
All internal control edges therefore remain identical.  Ordinary
fallthrough from $A$ enters $B$ at ordinal $a$, from $B$ enters $C$ at
ordinal $a+b$, and from $C$ falls through the full tree, in both layouts.
It remains to inspect escaping control edges.

For an instruction in $B$, its local prefix changes from $10$ to $01$.
An escaping break $B_k$ becomes $B_{k-1}$ in either layout, while an
escaping continue $C_1$ returns to ordinal zero in either layout and
$C_k$ for $k\geq2$ becomes $C_{k-1}$ in either layout.  Thus $B$ needs
no restriction.  For an instruction in $A$, the prefix changes from $0$
to $00$.  Escaping continues are unaffected, while escaping breaks would
change their targets; the hypothesis excludes precisely these breaks.
For an instruction in $C$, the prefix changes from $11$ to $1$.
Escaping breaks are unaffected, while escaping continues would change
their targets; the second hypothesis excludes precisely these continues.
Halting has no successor to compare.  These cases cover every instruction
outcome, proving the stated control-graph equality.\hfill$\square$

The two conditions are also necessary for equality of the complete local
control signatures with distinct symbolic external targets.  An escaping
break from $A$ with residual count one changes from $B_0$ to ordinal
$a+b$; with residual count $k\geq2$, it changes from $B_{k-1}$ to
$B_{k-2}$.  An escaping continue from $C$ with residual count one
changes from ordinal $a$ to ordinal zero; residual count two changes
from zero to $C_1$; and count $k\geq3$ changes from $C_{k-2}$ to
$C_{k-1}$.  This necessity concerns the complete syntactic control graph.
Unreachable instructions or special surrounding contexts can permit
further semantic identities.

\paragraph{Relation to the preceding rules.}
Requiring all three chunks to be syntactically closed is a sufficient
special case.  So is right association across a final break, since a
break leaf has no escaping continue.  The new criterion also permits
escaping continues in $A$, escaping breaks in $C$, arbitrary control exits
in $B$, and halt leaves.  A compiler search can test the two conditions
using just two nonnegative boundary counters per DAG node.

If $z(P)$ and $o(P)$ are the maximum numbers of zero and one ancestors
still needed outside $P$, then
\[
\begin{split}
z(\langle P,Q\rangle)&=\max\{\max(z(P)-1,0),z(Q)\},\\
o(\langle P,Q\rangle)&=\max\{o(P),\max(o(Q)-1,0)\}.
\end{split}
\]
For a decrement leaf these counters are $(1,0)$; for a break of depth
$k$, $(k,0)$; for a continue of depth $k$, $(0,k)$; and for an increment
or halt, $(0,0)$.  The admissibility test is exactly $z(A)=o(C)=0$.

\paragraph{Certificate discipline.}
The accompanying script independently computes local control signatures
from literal instruction paths and by composition of child signatures.
It checks their agreement on all live nodes of the inherited boundary
source and checks every initially admissible rotation against those
signatures.  A bounded beam search then uses the proved two-counter
criterion to explore source DAGs with the same instruction order and
control graph.  A separate checker verifies the complete final source
control graph, including both decrement branches.  These source
certificates do not establish a binary compiler theorem: changes in
source parentheses can disturb the inherited backward routing helpers
even when every register-machine instruction is preserved.

\paragraph{Concrete bounded-search result.}
The inherited boundary source has 6,455 leaves, 224 distinct source
subtrees, maximum instruction-path length 48, and maximum consecutive-zero
run six.  A beam of width 150 over 20 levels examined 625,858 generated
alternatives, including 40,486 whose chosen rotation was outside the
previous all-closed rule.  Counts refer to generated alternatives within
the beam search; alternatives can recur at different levels.

The best retained source has the same 6,455 leaves, the same 224 distinct
subtrees, and the same six-zero bound, but maximum instruction-path length
\textbf{39}.  Its replay records 19 transformations, five of them outside
the all-closed rule.  The independent full-graph checker verifies the
identity on every instruction and all 8,386 guarded control edges after
each recorded transformation.  The source record's SHA-256 is
\begin{quote}\small\ttfamily
b1cc93221d70b4ee28abc790aa032fcca461db9729e75910288147fd4aafbb02.
\end{quote}
The single-rotation checker independently examines all 328 possible local
rotations of the inherited source: the broader criterion admits 245,
compared with 225 admitted by the all-closed rule.  Every acceptance and
rejection agrees with the literal local control signature, and the
certificate includes concrete changed-target witnesses for rejections.

No 223-node source was found within the declared beam and caps.  The
39-depth source is a checked source-level alternative for later compiler
work; no transition table, machine state reduction, or runtime improvement
is claimed for it.  The achieved 238-state machine elsewhere in this
report continues to use the inherited boundary source.
